\documentclass[preview]{standalone}

\usepackage{amsmath}
\usepackage{amssymb}
\usepackage{bettelini}
\usepackage{stellar}
\usepackage{definitions}

\begin{document}

\id{measuretheory-convergence-theorems}
\genpage

\section{Convergence Theorems for Nonnegative Functions}

\begin{snippettheorem}{monotone-convergence-measurable-non-negative-function-theorem}{Monotone convergence theorem}
    Let \((X,\mathcal M,\mu)\) be a measure space and let
    \(f_n\colon X\fromto[0,+\infty]\) be measurable. If
    \(f_n\uparrow f\) pointwise, then
    \[
        \int_Xf_n\,d\mu\uparrow\int_Xf\,d\mu.
    \]
\end{snippettheorem}

The pointwise limit exists automatically because the sequence is increasing,
although it may take the value \(+\infty\).

\begin{snippetproof}{monotone-convergence-measurable-non-negative-function-theorem-proof}{monotone-convergence-measurable-non-negative-function-theorem}{Monotone convergence theorem}
    Monotonicity of the integral gives
    \[
        \int_Xf_n\,d\mu\leq\int_Xf_{n+1}\,d\mu,
    \]
    so these integrals increase to some \(\alpha\in[0,+\infty]\). Since
    \(f_n\leq f\), we have \(\alpha\leq\int_Xf\,d\mu\).

    For the reverse inequality, fix a nonnegative measurable simple function
    \(s\leq f\) and a constant \(0<c<1\). Define
    \[
        E_n=\{x\in X\suchthat f_n(x)\geq cs(x)\}.
    \]
    The sets \(E_n\) increase. Moreover, \(\bigcup_nE_n=X\): if \(f(x)=0\)
    then \(s(x)=0\), while if \(f(x)>0\), the strict inequality
    \(cs(x)<f(x)\) implies that \(f_n(x)\geq cs(x)\) eventually. Therefore
    \[
        \int_Xf_n\,d\mu
        \geq\int_{E_n}f_n\,d\mu
        \geq c\int_{E_n}s\,d\mu.
    \]
    The set function \(\varphi(E)=\int_Es\,d\mu\) is a measure, so continuity
    from below yields
    \[
        \int_{E_n}s\,d\mu=\varphi(E_n)\longrightarrow
        \varphi(X)=\int_Xs\,d\mu.
    \]
    Thus \(\alpha\geq c\int_Xs\,d\mu\). Letting \(c\uparrow1\) and then taking
    the supremum over all simple \(s\leq f\) gives
    \(\alpha\geq\int_Xf\,d\mu\).
\end{snippetproof}

\begin{snippetexample}{monotone-convergence-monotonicity-is-necessary}{Why monotonicity is necessary}
    On \(\realnumbers\) with Lebesgue measure, let
    \[
        f_n=n\chi_{(0,1/n)}.
    \]
    Then \(f_n\to0\) pointwise, but
    \[
        \int_\realnumbers f_n\,d\mu=1
        \neq0=\int_\realnumbers0\,d\mu.
    \]
    The sequence is not monotone.
\end{snippetexample}

\begin{snippetproposition}{nonnegative-lebesgue-integral-additivity}{Additivity of the nonnegative integral}
    If \(f,g\colon X\fromto[0,+\infty]\) are measurable, then
    \[
        \int_X(f+g)\,d\mu
        =\int_Xf\,d\mu+\int_Xg\,d\mu.
    \]
\end{snippetproposition}

\begin{snippetproof}{nonnegative-lebesgue-integral-additivity-proof}{nonnegative-lebesgue-integral-additivity}{Additivity of the nonnegative integral}
    Choose nonnegative measurable simple functions
    \(s_n\uparrow f\) and \(t_n\uparrow g\). Then
    \(s_n+t_n\uparrow f+g\), and additivity holds for simple functions.
    Applying monotone convergence to all three sequences gives
    \begin{align*}
        \int_X(f+g)\,d\mu
        &=\lim_n\int_X(s_n+t_n)\,d\mu\\
        &=\lim_n\left(\int_Xs_n\,d\mu+\int_Xt_n\,d\mu\right)\\
        &=\int_Xf\,d\mu+\int_Xg\,d\mu.
    \end{align*}
\end{snippetproof}

\begin{snippettheorem}{nonnegative-series-integration-theorem}{Integration of a nonnegative series}
    If \(f_n\colon X\fromto[0,+\infty]\) are measurable, then
    \[
        \int_X\sum_{n=1}^{\infty}f_n\,d\mu
        =\sum_{n=1}^{\infty}\int_Xf_n\,d\mu.
    \]
\end{snippettheorem}

\begin{snippetproof}{nonnegative-series-integration-theorem-proof}{nonnegative-series-integration-theorem}{Integration of a nonnegative series}
    By finite additivity of the integral,
    \[
        \int_X\sum_{n=1}^Nf_n\,d\mu
        =\sum_{n=1}^N\int_Xf_n\,d\mu.
    \]
    The partial sums increase pointwise to \(\sum_{n\geq1}f_n\), so the
    monotone convergence theorem gives the result as \(N\to\infty\).
\end{snippetproof}

\begin{snippettheorem}{fatou-lemma-theorem}{Fatou's lemma}
    Let \(f_n\colon X\fromto[0,+\infty]\) be measurable. Then
    \[
        \int_X\liminf_{n\to\infty}f_n\,d\mu
        \leq
        \liminf_{n\to\infty}\int_Xf_n\,d\mu.
    \]
\end{snippettheorem}

\begin{snippetproof}{fatou-lemma-theorem-proof}{fatou-lemma-theorem}{Fatou's lemma}
    Put
    \[
        g_k=\inf_{n\geq k}f_n.
    \]
    Each \(g_k\) is measurable, \(g_k\uparrow\liminf_nf_n\), and
    \(g_k\leq f_k\). Monotone convergence and monotonicity of the integral
    therefore give
    \[
        \int_X\liminf_nf_n\,d\mu
        =\lim_k\int_Xg_k\,d\mu
        \leq\liminf_k\int_Xf_k\,d\mu.
    \]
\end{snippetproof}

For nonpositive measurable functions the analogous statement uses
\(\limsup\) and the reverse inequality.

\section{Measures Defined by Densities}

Let \(f\colon X\fromto[0,+\infty]\) be measurable and define
\[
    \varphi_f(E)\triangleq\int_Ef\,d\mu.
\]

\begin{snippetproposition}{measure-defined-by-density}{A density defines a measure}
    The set function \(\varphi_f\) is a measure on \(\mathcal M\). Moreover,
    for every nonnegative measurable \(g\),
    \[
        \int_Xg\,d\varphi_f=\int_Xfg\,d\mu.
    \]
\end{snippetproposition}

\begin{snippetproof}{measure-defined-by-density-proof}{measure-defined-by-density}{A density defines a measure}
    If the sets \(E_i\) are pairwise disjoint, then
    \[
        \chi_{\bigcup_{i=1}^{\infty}E_i}
        =\sum_{i=1}^{\infty}\chi_{E_i}.
    \]
    Integration of a nonnegative series gives
    \begin{align*}
        \varphi_f\!\left(\bigcup_{i=1}^{\infty}E_i\right)
        &=\int_Xf\chi_{\bigcup_iE_i}\,d\mu\\
        &=\sum_{i=1}^{\infty}\int_Xf\chi_{E_i}\,d\mu
        =\sum_{i=1}^{\infty}\varphi_f(E_i).
    \end{align*}
    The change-of-measure identity is immediate for indicator functions,
    follows by linearity for nonnegative simple functions, and then follows
    for arbitrary nonnegative measurable functions by monotone convergence.
\end{snippetproof}

We write \(d\varphi_f=f\,d\mu\), or
\[
    f=\frac{d\varphi_f}{d\mu}.
\]
When such an \(f\) exists, it is called the Radon--Nikodym derivative of
\(\varphi_f\) with respect to \(\mu\).

\begin{snippetdefinition}{absolute-continuity-of-measures-definition}{Absolute continuity of measures}
    Let \(\varphi\) and \(\mu\) be measures on the same measurable space.
    The measure \(\varphi\) is \emph{absolutely continuous} with respect to
    \(\mu\), written \(\varphi\ll\mu\), if
    \[
        \mu(E)=0\implies\varphi(E)=0
    \]
    for every measurable set \(E\).
\end{snippetdefinition}

If \(d\varphi=f\,d\mu\), then necessarily \(\varphi\ll\mu\). The
Radon--Nikodym theorem gives a converse under suitable hypotheses.

\end{document}
